\section{Limitations}

The pilot is small: 60 latent tasks, two neutral metadata views, four model labels, and one prompt/schema condition. It is not statistically powered for fine-grained provider ranking. Reported model differences should be treated as descriptive.

The tasks are synthetic. Synthetic control is useful for isolating source-role failures, but it may introduce prose regularities that real evidence environments do not share. The leakage and baseline gates reduce this risk; they do not eliminate it.

The human review evidence has two different scopes. The 400-row pre-model surface review is a local leakage gate. The 40-row scorer audit is a local Codex-assisted manual audit. Neither should be described as independent human validation.

The Phase 1.1 schema-ablation slice is small: 16 source tasks, one visible metadata view, two main models, and a local Codex-assisted scorer audit. It supports only an interface-sensitivity claim about machine-readable evidence-role allocation. It should not be treated as a new benchmark result, a provider ranking, or evidence that schema wording explains every Phase 1 failure.

Finally, the artifact includes hidden labels and gold data for reproducibility, but those files are not model-visible. Any downstream reuse must preserve that separation.
